% Related work — IEEE BigData 2026 High School Symposium
% Target: ~0.5 column-page. Cite keys resolve against ../../refs.bib
%
% Discipline notes (docs/protocol_incomparability_note.md 4):
%   - descriptive, never accusatory; claim incomparability, never wrongness
%   - the spread across the field is the evidence; LFANN is the worked example
%   - state plainly that the random-init control is NOT ours

\section{Related Work}

\textbf{EEG foundation models.}
LaBraM~\cite{jiang2024labram} and CBraMod~\cite{wang2025cbramod} are
transformer encoders pretrained on large unlabelled EEG corpora and adapted
downstream by linear probing or fine-tuning. We audit both.

\textbf{Benchmarks, and their disagreement.}
At least seven recent efforts benchmark EEG foundation
models~\cite{wu2025adabrainbench,xiong2025eegfmbench,shen2026brain4fms,%
banville2026neuralbench,kontras2026neuroatlas,lu2026omnieegbench,%
zare2026negativecontrol}, and their verdicts do not agree. NeuroAtlas, spanning
42 datasets and 260k hours, concludes that pretrained models perform ``largely
on par'' with supervised baselines and that EEG-specific models do not
consistently outperform generic time-series
models~\cite{kontras2026neuroatlas}. Zare reports that a matched randomly
initialised encoder scored \emph{above} a pretrained encoder on a clinical
cognition task, and that paired confidence intervals leave comparator
superiority unresolved~\cite{zare2026negativecontrol}.

\textbf{The random-initialisation control is established methodology; we adopt
it rather than introduce it.}
Matched random-weight encoders are already run by
NeuroAtlas~\cite{kontras2026neuroatlas},
NeuralBench~\cite{banville2026neuralbench}, and
Zare~\cite{zare2026negativecontrol}. What differs here is the target. Those
audits evaluate predominantly clinical endpoints---cognition, seizure, sleep
staging, brain age---which differ from four-class motor imagery in trial count,
chance level, and the nature of the supervised comparator. Zare additionally
evaluates CBraMod on clinical cognition, where it trails classical features
(0.669 versus 0.734 macro-AUROC); we report the same direction on a different
task family, under a different protocol.

\textbf{Preprocessing: a named confound, not yet a measured one.}
NeuroAtlas observes that published evaluations differ in ``the EEG-specific
preprocessing that might influence reported
results''~\cite{kontras2026neuroatlas}. We measure that term directly, by
retraining each supervised comparator on the exact broadband arrays the
foundation models consume. The resulting preprocessing effect \emph{changes
sign across architectures}, so a matched-input control run on a single
architecture can support whichever attribution its author prefers. To our
knowledge this contrast has not previously been reported.

\textbf{Supervised motor-imagery decoders.}
Compact convolutional decoders such as
EEGNet~\cite{lawhern2018eegnet} remain widely used reference points. We compare
against three such decoders under one protocol:
ShallowConvNet~\cite{schirrmeister2017deep}, ATCNet~\cite{altaheri2023atcnet},
and EEG Conformer~\cite{song2023conformer}.

\textbf{Evaluation protocol.}
Published four-class BCI~IV-2a accuracies span roughly 68\% to 91\%, and the
spread tracks evaluation protocol at least as much as architecture. Reported
choices include augmentation applied before partitioning, within-session random
splits in place of the official cross-session protocol, and unstated checkpoint
selection~\cite{yu2022lffn}. Such numbers are not commensurable with
cross-session results, and are frequently tabulated as though they were.
MOABB~\cite{jayaram2018moabb} addresses this by standardising
pipelines; we address it by locking every selection decision to a
training-derived validation split and enforcing that lock mechanically.

\textbf{Calibration.}
We report calibration alongside accuracy, using temperature
scaling~\cite{guo2017calibration} fitted on validation data only, and proper
scoring rules~\cite{brier1950verification,gneiting2007proper}.
